\documentclass[aspectratio=169,10pt]{beamer}
\usetheme[numbering=fraction,progressbar=frametitle]{metropolis}
\usepackage{fontspec}
\setsansfont{Latin Modern Sans}
\usepackage{amsmath,amssymb,mathtools,graphicx,booktabs}
\usepackage{polyglossia}
\setdefaultlanguage{english}
\setbeamertemplate{navigation symbols}{}
\setbeamerfont{frametitle}{size=\large}
\setbeamerfont{normal text}{size=\normalsize}
\setlength{\parskip}{.45em}
\hypersetup{pdfauthor={},pdfsubject={Lecture slides},pdfkeywords={},pdftitle={Differentiation: rates and tangents}}
\providecommand{\N}{\mathbb{N}}
\providecommand{\Z}{\mathbb{Z}}
\providecommand{\D}{\mathbb{D}}
\providecommand{\Q}{\mathbb{Q}}
\providecommand{\R}{\mathbb{R}}
\providecommand{\C}{\mathbb{C}}
\providecommand{\vect}[1]{\overrightarrow{#1}}
\providecommand{\vi}{\vec{\imath}}
\providecommand{\vj}{\vec{\jmath}}
\providecommand{\vk}{\vec{k}}
\providecommand{\norm}[1]{\left\lVert #1\right\rVert}
\providecommand{\vc}[2]{\begin{pmatrix}#1\\#2\end{pmatrix}}
\providecommand{\deter}[2]{\det\left(#1,#2\right)}
\providecommand{\abs}[1]{\left\lvert #1\right\rvert}
\providecommand{\dd}{\,\mathrm{d}}
\providecommand{\e}{\mathrm{e}}
\providecommand{\Card}{\operatorname{Card}}
\providecommand{\E}{\operatorname{E}}
\providecommand{\V}{\operatorname{V}}
\providecommand{\degre}{^{\circ}}
\providecommand{\un}[1]{\,\mathrm{#1}}
\providecommand{\sci}[2]{#1\times10^{#2}}
\providecommand{\pgcd}{\operatorname{PGCD}}
\providecommand{\ppcm}{\operatorname{PPCM}}
\title{Differentiation: rates and tangents}
\subtitle{Lecture and worked examples}
\date{}
\author{}
\begin{document}
\begin{frame}[plain]
\vfill
{\usebeamerfont{title}\inserttitle\par}
\medskip
{\usebeamerfont{subtitle}\insertsubtitle\par}
\medskip
{\usebeamercolor[fg]{progress bar}\rule{\textwidth}{.5pt}}
\vfill
\end{frame}
\begin{frame}{Chapter path}
\textbf{Objectives.} Interpret a rate; approximate a derivative at a point; determine a tangent; understand why differentiability fails.

\textbf{Prerequisites.} Function values, slope of a line, algebraic manipulation, square root and absolute value.

\textbf{Yearly reference point.} 3-12 November: 8 hours
\end{frame}
\begin{frame}{Rate of change and secants}
For \(a\ne b\) in the domain of f, the rate of change between a and b is \(\dfrac{f(b)-f(a)}{b-a}\). It is the slope of the secant through the corresponding points. In context, the unit of the rate is the unit of f divided by that of the variable.
\end{frame}
\begin{frame}{Rate of change and secants}
For \(f(x)=x^2\), between 1 and 3, the rate is \(\dfrac{9-1}{3-1}=4\). Between 1 and \(1+h\), it is \(2+h\). The slopes for \(h=0.1\) and \(h=-0.1\) are 2.1 and 1.9.
\end{frame}
\begin{frame}{Rate of change and secants}
The rate is not \(\dfrac{f(b)}{f(a)}\); h represents a gap, positive or negative depending on the domain.
\end{frame}
\begin{frame}{Rate of change and secants}
Calculate the rate of \(f(x)=x^2+3x\) between 2 and \(2+h\).
\end{frame}
\begin{frame}{Derivative at a point}
f is differentiable at a when the rate \(\dfrac{[f(a+h)-f(a)]}{h}\) tends to a real number as h tends to 0, within the domain. This real number is \(f\prime(a)\). An intuitive approach to the limit is used here. In kinematics, if f gives a position in metres and time is in seconds, the derivative at a point represents an instantaneous velocity in \(\mathrm{m\,s^{-1}}\).
\end{frame}
\begin{frame}{Derivative at a point}
For \(s(t)=t^2+3t\), the rate around 2 is \(7+h\), so \(s\prime(2)=7\ \mathrm{m\,s^{-1}}\). The average velocity on \([2,3]\) is \(8\ \mathrm{m\,s^{-1}}\); it is not the velocity at time 2.
\end{frame}
\begin{frame}{Derivative at a point}
We do not replace h by 0 in an unsimplified quotient: we would get a division by 0.
\end{frame}
\begin{frame}{Derivative at a point}
For \(f(x)=3x^2-2x\), calculate \(f\prime(1)\) from the rate.
\end{frame}
\begin{frame}{Tangent and affine approximation}
If f is differentiable at a, the tangent at the point \((a,f(a))\) has equation \(y=f(a)+f\prime(a)(x-a)\). For small h, we use \(f(a+h)\approx f(a)+f\prime(a)h\). This local approximation does not by itself provide a bound on the error.
\end{frame}
\begin{frame}{Tangent and affine approximation}
\begin{columns}[T,onlytextwidth]
\begin{column}{.57\textwidth}
For \(f(x)=x^2\) at 2, tangent \(y=4+4(x-2)=4x-4\). We approximate \(2.01^2\) by \(4+4(0.01)=4.04\); the exact value is \(4.0401\), error \(0.0001=h^2\).
\end{column}
\begin{column}{.4\textwidth}
\centering
\includegraphics[width=\linewidth,height=4.5cm,keepaspectratio]{figures/premiere-04-1-en.pdf}
\end{column}
\end{columns}
\end{frame}
\begin{frame}{Tangent and affine approximation}
A tangent may cross the curve again elsewhere. “Only one common point” is not a definition.
\end{frame}
\begin{frame}{Tangent and affine approximation}
Given \(f(3)=5\) and \(f\prime(3)=-2\). Write the tangent and approximate \(f(3.04)\).
\end{frame}
\begin{frame}{When differentiability fails}
The square root function is defined on \([0,+\infty)\) but is not differentiable at 0. The absolute value function is defined on \(\R\) but is not differentiable at 0. Being defined, or having a continuous curve, therefore does not guarantee differentiability.
\end{frame}
\begin{frame}{When differentiability fails}
For \(|x|\), the secants from the origin all have slope 1 on the right and \(-1\) on the left. The corner point does not have a unique tangent with finite slope. For the square root, the limiting direction is vertical; there is no real derivative.
\end{frame}
\begin{frame}{When differentiability fails}
A very large slope is not a proof of non-differentiability; the behaviour of the rate must be explained.
\end{frame}
\begin{frame}{When differentiability fails}
Justify that \(g(x)=|x-2|\) is not differentiable at 2. Also study the differentiability of \(x\mapsto\sqrt{x}\) at 0 through its rate.
\end{frame}
\end{document}
